\documentclass{article}

\usepackage{float}
\restylefloat{table}

\usepackage{booktabs}

\title{Team Contributions: POC\\\progname}

\author{\authname}

\date{}

\input{../Comments}
\input{../Common}

\begin{document}

\maketitle

\noindent This document summarizes the contributions of each team member up to the POC Demo. The time period of interest is the time between the beginning of the term
and the POC demo.

\section{Demo Plans}

For the POC demo we plan to showcase the two major risks of our project which were identified in the \href{https://github.com/M9Huynh/technically-functional/blob/main/docs/DevelopmentPlan/DevelopmentPlan.pdf}{development plan document}. A short section has been copied for reference. \\

\noindent A few major risks in our proof of concept demonstration are:
\begin{itemize}
  \item Risks associated with measuring a 3D motion using a 2D view/plane. Currently the scope of the project includes patients performing an exercise of the knee as other joints add complexity with their degrees of freedom (i.e. shoulder, ankle, etc.).
  \begin{itemize}
    \item To address this, the team will test whether a  single camera view is sufficient for recording the movement. This will be tested primarily on the sagittal view and compared with adding frontal and transverse views. The accuracy of video analysis will be measured by comparing the difference in guided prompts and verifying if they can adhere to the requirements of the exercise (i.e. does a single view yield false positives? Does having two views change the accuracy achieved?).
  \end{itemize}  
  \item The second is being able to measure metrics related to stretching and past history. This was identified as a key risk since range of motion is a key metric in physiotherapy and measuring it over time is a good indicator of exercise efficacy. Furthermore, while talking with our supervisor, it was apparent that patients may not “push themselves” due to the perceived fragility of their post-operative condition.
  \begin{itemize}
    \item The system tracks a patient's rehabilitation progress by measuring their range of motion (i.e. the maximum/minimum angle of knee flexion) and monitoring exercise duration. Thus, the success criteria is if the team can identify range of motion and exercise duration.
  \end{itemize}
\end{itemize}


\section{Team Meeting Attendance}

\begin{table}[H]
\centering
\begin{tabular}{ll}
\toprule
\textbf{Student} & \textbf{Meetings}\\
\midrule
Total & 15\\
Matthew & 15\\
Cieran & 13\\
Vaisnavi & 14\\
Maham & 14\\
Eman & 14\\
\bottomrule
\end{tabular}
\end{table}


\noindent The above table does not include attendance for TA meetings during tutorial as well as smaller meetings for 'sub team work' after the VnV since not all team member's attendance was needed.

\section{Supervisor/Stakeholder Meeting Attendance}


\noindent \textbf{Supervisor's Name: Kevin Yu} 

\begin{table}[H]
\centering
\begin{tabular}{ll}
\toprule
\textbf{Student} & \textbf{Meetings}\\
\midrule
Total & 2\\
Matthew & 2\\
Cieran & 2\\
Vaisnavi & 2\\
Maham & 2\\
Eman & 2\\
\bottomrule
\end{tabular}
\end{table}

\noindent Currently the total number of (virtually) meetings have been 2 but we have periodically provided Kevin with updates on our project and we are progressing.

\section{Lecture Attendance}


\begin{table}[H]
\centering
\begin{tabular}{ll}
\toprule
\textbf{Student} & \textbf{Lectures}\\
\midrule
Total & 7\\
Matthew & 7\\
Cieran & 4\\
Vaisnavi & 4\\
Maham & 6\\
Eman & 2\\
\bottomrule
\end{tabular}
\end{table}

\noindent This table takes into account the software lectures from our initial group forming on Sep 15. This also does not take into account the guest lecture on Sep 29.

\section{TA Document Discussion Attendance}


\noindent \textbf{TA's Name: Rashad Bhuiyan}

\begin{table}[H]
\centering
\begin{tabular}{ll}
\toprule
\textbf{Student} & \textbf{Meetings}\\
\midrule
Total & 3\\
Matthew & 3\\
Cieran & 3\\
Vaisnavi & 3\\
Maham & 3\\
Eman & 2\\
\bottomrule
\end{tabular}
\end{table}

\noindent These meeting were for the Problem Statement and Development Plan, SRS and HA, and the VnV Plan.

\section{Commits}


\begin{table}[H]
\centering
\begin{tabular}{lll}
\toprule
\textbf{Student} & \textbf{Commits} & \textbf{Percent}\\
\midrule
Total & 190 & 100\% \\
Matthew & 47 & 24.7\% \\
Cieran & 53 & 27.9\% \\
Vaisnavi & 40 & 21.1\% \\
Maham & 46 & 24.2\% \\
Eman & 4 & 2.1\% \\
\bottomrule
\end{tabular}
\end{table}

Regarding Vaisnavi's commit count, she found out that her GitHub account was not properly linked to her local Git configuration at the time. As a result, many of her initial commits that were pushed were not linked under her name. After realizing this following the completion of the first deliverables (the two documents), she immediately corrected her Git configuration so that all of her following commits would be properly authored under her account. It was too late to co-author the commits, but she did leave a comment under each of the commits for those previous deliverables of the ones that she contributed to for her sections. It was too late to be added to the official commit count however at that point, but just wanted to point that out. As for Eman, she had not accepted the repo invite from Matthew and it ended up expiring, and as a result, could not commit until the end of the first deliverable. Additionally, Eman also has a preference for doing her work on other platforms (Google docs) and adds her commits at the very end. Therefore, her commit number is not reflective of her contribution towards the project. After some discussions and to make deliverable submissions smoother in the future, we concurrently decided to commit more often and make additional pull requests every 2 days (approximately). This will ensure that merge conflicts are kept to a minimum because they are time consuming to fix. 

\section{Issue Tracker}


\begin{table}[H]
\centering
\begin{tabular}{lll}
\toprule
\textbf{Student} & \textbf{Authored (O+C)} & \textbf{Assigned (C only)}\\
\midrule
Matthew & 73 & 13 \\
Cieran & 9 & 5 \\
Vaisnavi & 2 & 8 \\
Maham & 5 & 14 \\
Eman & 0 & 4 \\
\bottomrule
\end{tabular}
\end{table}

A majority of the assigned column is not done properly and mainly encapsulates the SRS and HA sections as that was the only deliverable that created issues for each section and were assigned accordingly but that was based on a different SRS template and does not capture the \href{https://github.com/M9Huynh/technically-functional/blob/main/docs/SRS/SRS.pdf}{template that we used} (Volere).

\section{CICD}

Currently, CI/CD’s use in the project is limited to LaTeX file compilation and has been for the deliverables submitted so far. As the proof of concept takes a more concrete shape within the coming weeks, as well as Rev. 0 and Rev. 1 in the coming months, CI/CD will be used more effectively on code compilation checks and general test cases for input fields, basic actions, and expected outputs.

While it’s currently uncertain how to integrate CI/CD with the main pain point of our project; the video submission feature which analyzes and provides live feedback on user form, there will be some experimenting and testing of inputs and outputs for the video analysis module. Currently, an accepted range of the video analyzed is expected to be the solution. This way, if the point analyzed by the model slightly shifts there is an acceptable range within a couple degrees to allow differences.

CI/CD used in this way will ensure the project is moving forward without leaving broken pieces behind and can continually assure we’re meeting important points, such as requirements testable upon each pull request.

\section{Team Charter Trigger Items}

\wss{Provide a summary of the quantified triggers identified in the team's
charter.}

\wss{Provide a list of any violations of the triggers.  If the team wishes, the
violations can be summarized on aggregate, instead of naming specific team
members.}

\wss{Provide a plan to address the violations.  This could include revising the
triggers, if they are found to be too weak, strong or ambiguous.}

\section{Additional Productivity Metrics}

\subsection{Summary of Quantified Triggers}
\begin{itemize}
    \item Missing 3 consecutive meetings when deliverables are not the main discussion topic triggers a team meeting
    \item Missing 2 consecutive meetings when deliverables are the main topic triggers a team meeting
    \item Deliverables must be finalized 2 days prior to the deadline (internal deadline)
\end{itemize}
 
\subsection{Violations}
No violations have been recorded at this time.
 
\subsection{Plan to Address Violations}
If triggers are violated, the team will:

\begin{enumerate}
    \item Hold a team meeting to discuss the attendance or deadline issues
    \item Review the circumstances with the member involved
    \item For extended emergencies ($>$2 weeks), bring the matter to teaching staff to ensure workload remains manageable for the team
    \item If triggers are found to be too weak, strong, or ambiguous, they will be revised through the team's decision-making process (this will be done through a unanimous vote or majority vote with all members heard)
\end{enumerate}

\textit{Duration between pull requests:} Over the course of the project, our team learned that more frequent pull requests would be beneficial to the workflow. Frequent merges will ensure minimal merge conflicts towards the end of the deliverable and guarantee that any output is up-to-date. Merge conflicts are time-consuming and troublesome to handle, and implementing this strategy will result in an overall improvement in the quality of our work. 


\end{document}